%--------------------
% Packages
% -------------------
\documentclass[11pt,a4paper]{article}
\usepackage[utf8x]{inputenc}
\usepackage[T1]{fontenc}
\usepackage{amsmath, amssymb, amsthm}
\usepackage{algorithm, algpseudocode}

\usepackage[pdftex]{graphicx} % Required for including pictures
\usepackage[swedish]{babel} % Swedish translations
\usepackage[pdftex,linkcolor=black,pdfborder={0 0 0}]{hyperref} % Format links for pdf
\usepackage{calc} % To reset the counter in the document after title page
\usepackage{enumitem} % Includes lists
\usepackage{xcolor} % Includes colour

\frenchspacing % No double spacing between sentences
\linespread{1.2} % Set linespace
\usepackage[a4paper, lmargin=0.1666\paperwidth, rmargin=0.1666\paperwidth, tmargin=0.1111\paperheight, bmargin=0.1111\paperheight]{geometry} %margins
%\usepackage{parskip}


%--------------------
% Title page
% -------------------
\title{Final report: L-BFGS Methods for Periodic Orbit Search}
\author{Maximilian Hartenberger}
\date{\today}

%--------------------
% Document
% -------------------
\begin{document}
\maketitle

\begin{abstract}
% TODO: Write the abstract here.
\end{abstract}

\section{Introduction}
% TODO: Write the introduction here. Blablah about periodic orbits lying dense in an attractor and the importance of finding them. Mention the L-BFGS and L-BFGS-Dogleg methods and their relevance to the problem.

\section{Problem Formulation}
\subsection{Periodic orbits}
A periodic orbit of period $T>0$ is a solution of the initial value problem \eqref{eq:ivp} such that $\varphi(T,x_0) = x_0$.

\subsection{Residual and Derivative}
Let \begin{equation*}
    z = \begin{pmatrix} u_1 \\ \vdots \\ u_N \\ T \end{pmatrix}
\end{equation*}
where $u_1, \dots, u_N$ form a partition of the periodic orbit and $T$ is the period. We have $\varphi(T/N,u_i) = u_{i+1}$ for $i=1,\dots,N-1$ and $\varphi(T/N,u_N) = u_1$. Thus, we can define the residual
\begin{align*}
    \tilde{F}: \mathbb{R}^{nN+1} &\to \mathbb{R}^{nN} \\
    z &\mapsto \begin{pmatrix} \varphi(T/N,u_1)-u_2 \\ \vdots \\ \varphi(T/N,u_{N-1})-u_N \\ \varphi(T/N,u_N)-u_1 \end{pmatrix}.
\end{align*}
The periodic orbit search can be formulated as the root finding problem $\tilde{F}(z) = 0$. The derivative of $\tilde{F}$ is given by \begin{equation}
    D\tilde{F}(z) = \begin{pmatrix} \Phi(T/N,u_1) & -I & 0 & \dots & 0 & \frac{1}{N}f(\varphi(T/N,u_1)) \\ 0 & \Phi(T/N,u_2) & -I & \dots & 0 & \frac{1}{N}f(\varphi(T/N,u_2)) \\ \vdots & \vdots & \vdots & \ddots & \vdots & \vdots \\ -I & 0 & 0 & \dots & \Phi(T/N,u_N) & \frac{1}{N}f(\varphi(T/N,u_N))\end{pmatrix}.
\end{equation}

\subsection{Residual and Derivative with spatial shift}

For dynamical systems with continuous symmetries (e.g., translation or rotation invariance), a periodic orbit may close only up to a symmetry transformation — a so-called \textit{relative periodic orbit}. Let
\begin{equation*}
    S: \mathbb{R}^n \times \mathbb{R} \to \mathbb{R}^n, \quad (u, s) \mapsto S(u, s)
\end{equation*}
be a spatial shift operator with scalar shift parameter $s$. The state vector is extended to
\begin{equation*}
    z = \begin{pmatrix} u_1 \\ \vdots \\ u_N \\ T \\ s \end{pmatrix} \in \mathbb{R}^{nN+2},
\end{equation*}
and the closing condition on the last segment becomes
\begin{equation*}
    S\bigl(\varphi(T/N, u_N),\, s\bigr) = u_1.
\end{equation*}
The unaugmented residual $\tilde{F}: \mathbb{R}^{nN+2} \to \mathbb{R}^{nN}$ is then
\begin{align*}
    \tilde{F}_i(z) &= \varphi(T/N, u_i) - u_{i+1}, \qquad i = 1, \dots, N-1, \\[2pt]
    \tilde{F}_N(z) &= S\bigl(\varphi(T/N, u_N),\, s\bigr) - u_1.
\end{align*}
Let $\varphi_i := \varphi(T/N, u_i)$ and $\Phi_i := \Phi(T/N, u_i)$.  Defining the shifted quantities
\begin{equation*}
    \widehat{\Phi}_N := \partial_u S(\varphi_N, s)\,\Phi_N, \qquad
    \widehat{f}_N   := \partial_u S(\varphi_N, s)\,f(\varphi_N),
\end{equation*}
the derivative $D\tilde{F}(z) \in \mathbb{R}^{nN \times (nN+2)}$ is
\begin{equation}\label{eq:jac_spatial_shift}
    D\tilde{F}(z) = 
    \begin{pmatrix}
        \Phi_1      & -I          & 0           & \dots  & 0              & \frac{1}{N}f(\varphi_1)       & 0 \\[4pt]
        0           & \Phi_2      & -I          & \dots  & 0              & \frac{1}{N}f(\varphi_2)       & 0 \\[4pt]
        \vdots      & \vdots      & \vdots      & \ddots & \vdots         & \vdots                        & \vdots \\[4pt]
        -I          & 0           & 0           & \dots  & \widehat{\Phi}_N & \frac{1}{N}\widehat{f}_N   & \partial_s S(\varphi_N, s)
    \end{pmatrix}.
\end{equation}

\subsubsection*{Phase conditions}
For a relative periodic orbit we need \emph{two} phase conditions to remove the degeneracies from time translation and spatial shift:
\begin{align*}
    \langle u_1 - u_{\text{ref}},\; f(u_{\text{ref}}) \rangle &= 0 \quad &\text{(time phase)}, \\[2pt]
    \langle u_1 - u_{\text{ref}},\; g(u_{\text{ref}}) \rangle &= 0 \quad &\text{(spatial phase)},
\end{align*}
where $g(u) := \partial_s S(u, 0)$ is the infinitesimal generator of the symmetry.  Choosing $u_{\text{ref}} = u_1$ makes both residuals identically zero.  Appending the two phase-condition rows to~\eqref{eq:jac_spatial_shift} yields the augmented, square Jacobian $J(z) \in \mathbb{R}^{(nN+2)\times (nN+2)}$:
\begin{equation*}
    J(z) = \left(\begin{array}{c}
        \multicolumn{7}{c}{\raisebox{0.8ex}{\mbox{\footnotesize $D\tilde{F}(z)$}}} \\[2pt]
        f(u_1)^\top & 0 & 0 & \dots & 0 & 0 & 0 \\[2pt]
        g(u_1)^\top & 0 & 0 & \dots & 0 & 0 & 0
    \end{array}\right).
\end{equation*}
The corresponding augmented residual is
\begin{equation*}
    F(z) = \begin{pmatrix}
        \varphi(T/N, u_1) - u_2 \\
        \vdots \\
        \varphi(T/N, u_{N-1}) - u_N \\
        S(\varphi(T/N, u_N), s) - u_1 \\
        0 \\ 0
    \end{pmatrix} \in \mathbb{R}^{nN+2}.
\end{equation*}

In the implementation, the operator $S$ is passed as an argument (\texttt{S} or \texttt{nothing} when no spatial symmetry is present).  The forward matrix-vector product $J(z)\cdot\delta z$ and its transpose $J(z)^\top \cdot w$ are computed matrix-free in \texttt{Stage\-Iter\-Cache} and \texttt{AdjointIterSolCache}, respectively.  For the adjoint, the inverse shift $S(\cdot, -s)$ is applied to the last segment \emph{before} the backward adjoint integration, matching the forward composition $S \circ D\varphi_N$.
The root finding problem $\tilde{F}(z) = 0$ is underdetermined since $D\tilde{F}(z)\in\mathbb{R}^{nN\times (nN+1)}$ is not full rank. To make it well-posed, we need to add a phase condition \begin{equation*}
    \langle u_1 - u_{ref}, f(u_{ref}) \rangle = 0
\end{equation*}
with derivative
\begin{equation*}
    \begin{pmatrix} f(u_{ref})^\top & 0 & \dots & 0 \end{pmatrix}.
\end{equation*} 
Adding the phase condition to the derivative gives the following augmented derivative
\begin{equation*}
    J(z)=\begin{pmatrix} \Phi(T/N,u_1) & -I & 0 & \dots & 0 & \frac{1}{N}f(\varphi(T/N,u_1)) \\ 0 & \Phi(T/N,u_2) & -I & \dots & 0 & \frac{1}{N}f(\varphi(T/N,u_2)) \\ \vdots & \vdots & \vdots & \ddots & \vdots & \vdots \\ -I & 0 & 0 & \dots & \Phi(T/N,u_N) & \frac{1}{N}f(\varphi(T/N,u_N))\\
        f(u_{ref})^T & 0 & 0 & \dots & 0 & 0 \end{pmatrix}.
\end{equation*}
where we choose $u_{ref} = u_1$.
Because of this the phase condition in the residual is always $0$ and we only need to add a $0$ to the residual to match the dimension and get the augmented residual
\begin{equation*}
    F(z) = \begin{pmatrix} \varphi(T/N,u_1)-u_2 \\ \vdots \\ \varphi(T/N,u_{N-1})-u_N \\ \varphi(T/N,u_N)-u_1 \\ 0 \end{pmatrix}.
\end{equation*}

\subsection{Root finding}
\subsubsection*{Classical root finding}
There exist two approaches to solve the root finding problem $F(z) = 0$: The first approach is to use Newton's method (with a globalization strategy) which requires the solution of a linear system
\begin{equation*}
    J(z_k)\delta z = -F(z_k)
\end{equation*}
and an update $z_{k+1} = z_k + \lambda\delta z$ at each iteration with a suitable step size $\lambda$.

\subsubsection*{Optimization problem}
The second approach is to formulate the root finding problem as an optimization problem
\begin{equation*}
    \min_{z\in\mathbb{R}^{nN+1}} \phi(z) = \frac{1}{2}\|F(z)\|^2.
\end{equation*}
Since it is an unconstrained problem, we get the solution by finding a root of the gradient $\nabla \phi(z) = J(z)^\dagger F(z)$.


\section{Mathematical Background}
\subsection{Nonlinear flow}
Let $f\in C^k(\mathbb{R}^n)$ and 
\begin{equation*}
    \varphi: \mathbb{R} \times \mathbb{R}^n \to \mathbb{R}^n, \quad (t,u_0) \mapsto u(t)
\end{equation*}
be the flow map where $u$ is the solution of the initial value problem
\begin{equation}\label{eq:ivp}
    \dot{u} = f(u), \quad u(0) = u_0.
\end{equation}
We have that \begin{equation}
    \frac{\partial \varphi}{\partial t}(t,u_0) = f(\varphi(t,u_0)).
\end{equation}
\subsection{Linearized flow}
Differentiating with respect to $u_0$ yields
\begin{equation*}
     \frac{\partial}{\partial u_0} \frac{\partial \varphi}{\partial t}(t,u_0) = \frac{\partial}{\partial u_0} f(\varphi(t,u_0)).
\end{equation*}
Swapping the order of differentiation and applying the chain rule gives
\begin{equation*}
    \frac{\partial}{\partial t} \frac{\partial \varphi}{\partial u_0}(t,u_0) = Df(\varphi(t,u_0))\frac{\partial \varphi}{\partial u_0}(t,u_0).
\end{equation*}
Denote $\Phi(t,u_0) := \frac{\partial \varphi}{\partial u_0}(t,u_0)$, then $\Phi$ is the solution of the initial value problem
\begin{equation}\label{eq:lin_flow}
    \frac{\partial \Phi}{\partial t} = Df(\varphi(t,u_0))\Phi, \quad \Phi(0,u_0) = I.
\end{equation}
$\Phi$ is called the \textit{linearized flow} of $\varphi$. \begin{align*}
    \Phi(\cdot,u_0):\mathbb{R}\times \mathbb{R}^n &\to \mathbb{R}^{n} \\
    (t,v_0) &\mapsto \Phi(t,u_0)\cdot v_0
\end{align*}
$v_0$ is called the \textit{initial variation} and $\Phi(t,u_0)\cdot v_0$ the \textit{variation} at time $t$. $v$ is the solution of the initial value problem
\begin{equation}\label{eq:lin_var}
    \dot{v} = Df(\varphi(t,u_0))v, \quad v(0) = v_0.
\end{equation}
For better understanding, we can look at the linearized flow as a linear approximation of the flow $\varphi$ in the following way: Let $v_0=u_0-\bar{u}_0$ be a small perturbation of the initial condition $\bar{u}_0$. Then using Taylor expansion we have
\begin{equation*}
    \varphi(t,\bar{u}_0) = \varphi(t,u_0) - \Phi(t,u_0)v_0 + \mathcal{O}(\|v_0\|^2).
\end{equation*}
The linearized flow $\Phi$ is the first order approximation of the flow $\varphi$ with respect to the initial condition.
\begin{align*}
    \Phi(t,u_0)(u_0-\bar{u}_0) &= \varphi(t,u_0) - \varphi(t,\bar{u}_0) + \mathcal{O}(\|u_0-\bar{u}_0\|^2)\\
    &\approx \varphi(t,u_0) - \varphi(t,\bar{u}_0).
\end{align*}
\subsection{Adjoint flow}
The Adjoint flow $\Psi$ is defined as the transpose of the linearized flow $\Phi$:
\begin{equation*}
    \Psi(t,u_0) := \Phi(t,u_0)^\top.
\end{equation*}
It maps the final variation at time $t$ to the initial variation at time $0$. Let $q_T$ be the final variation at time $T$. We have \begin{equation*}
    \Psi(T,u_0)q_T = q(0).
\end{equation*}
Define $q(t) := \Psi(T-t,\varphi(t,u_0))q_T$.
Using the cocycle property of the linearized flow $\Phi$:
\begin{equation*}
    \Phi(t+s,u_0) = \Phi(t,\varphi(s,u_0))\Phi(s,u_0)
\end{equation*}
We have
\begin{equation*}
    \Phi(T,u_0) = \Phi(T-t,\varphi(t,u_0))\Phi(t,u_0).
\end{equation*}
Taking the transpose gives the adjoint cocycle property \begin{equation*}
    \Psi(T,u_0) = \Psi(t,u_0)\Psi(T-t,\varphi(t,u_0)).
\end{equation*}

We get 
\begin{align*}
    q(0) &= \Psi(T,u_0)q_T \\
         &= \Psi(t,u_0)\Psi(T-t,\varphi(t,u_0))q_T \\
         &= \Psi(t,u_0)q(t).
\end{align*}
and thus
\begin{align*}
     q(t) &= \Psi(t,u_0)^{-1}q(0) \\
          &= \Phi(t,u_0)^{-\top}q(0)
\end{align*}
Differentiating with respect to $t$ and using the equality $\frac{\dot{M^{-1}}}{dt} = -M^{-1}\dot{M}M^{-1}$ for an invertible matrix $M$ gives
\begin{align*}
    \frac{\partial q}{\partial t}(t) &= -\Phi(t,u_0)^{-\top}\frac{\partial \Phi}{\partial t}(t,u_0)^{\top}\Phi(t,u_0)^{-\top}q(0) \\
    &= -\Phi(t,u_0)^{-\top}\left(Df(\varphi(t,u_0))\Phi(t,u_0)\right)^{\top}\Phi(t,u_0)^{-\top}q(0) \\
    &= -\Phi(t,u_0)^{-\top}\Phi(t,u_0)^{\top}Df(\varphi(t,u_0))^{\top}\Phi(t,u_0)^{-\top}q(0) \\
    &= -Df(\varphi(t,u_0))^{\top}q(t).
\end{align*}
Thus, the adjoint flow $\Psi$ is the solution of the initial value problem
\begin{equation}\label{eq:adjoint_flow}
    \frac{\partial \Psi}{\partial t} = -Df(\varphi(t,u_0))^\top\Psi, \quad \Psi(0,u_0) = I.
\end{equation} and the variation $q$ is the solution of the initial value problem
\begin{equation}\label{eq:variation}
    \dot{q} = -Df(\varphi(t,u_0))^\top q, \quad q(T) = q_T.
\end{equation}
We have the following duality relation between the linearized flow $\Phi$ and the adjoint flow $\Psi$: Let $v_0$ be an initial variation and $q_T$ a final variation. Then
\begin{align*}
    \langle q_T, v(T) \rangle &= \langle q_T, \Phi(T,u_0)v_0 \rangle \\
    &= \langle \Phi(T,u_0)^\top q_T, v_0 \rangle \\
    &= \langle \Psi(T,u_0) q_T, v_0 \rangle \\
    &= \langle q(0), v_0 \rangle.
\end{align*}
That is, the scalar product of the forward variation $v$ and the adjoint variation $q$ is conserved in time.

\section{Adjoint flow with spatial shift}

When a spatial shift operator $S$ is present, the Jacobian $J(z)$ on the last segment
involves the composition $S \circ D\varphi_N$ rather than $D\varphi_N$ alone.  The adjoint
(transpose) $J(z)^\top$ must therefore incorporate the transpose of this composition.
We derive the full action of $J(z)^\top \cdot w$ component by component.

\subsection{Forward action (recap)}

Recall from~\eqref{eq:jac_spatial_shift} that the Jacobian acts on a perturbation
$\delta z = (\delta u_1, \dots, \delta u_N, \delta T, \delta s)$ as
\begin{align}
    (J\,\delta z)_i &= -\Phi_i\,\delta u_i + \delta u_{i+1}
                       - \tfrac{1}{N}f(\varphi_i)\,\delta T,
                       \qquad i = 1,\dots,N-1, \label{eq:jfw_inner} \\[4pt]
    (J\,\delta z)_N &= -\widehat{\Phi}_N\,\delta u_N + \delta u_1
                       - \tfrac{1}{N}\widehat{f}_N\,\delta T
                       + \partial_s S(\varphi_N,s)\,\delta s, \label{eq:jfw_last}
\end{align}
where $\widehat{\Phi}_N = \partial_u S(\varphi_N,s)\,\Phi_N$ and
$\widehat{f}_N = \partial_u S(\varphi_N,s)\,f(\varphi_N)$.  The two phase-condition
rows (time and spatial) contribute only to the scalar components and are treated
separately.

\subsection{Adjoint of the spatial shift operator}

Write $\widehat{S}(\cdot) := \partial_u S(\varphi_N, s)$ for the Jacobian of $S$
with respect to its first (state) argument, evaluated at the current orbit point.
For a pure translation $S(u,s) = u + s\cdot\mathbf{e}$ with a fixed direction
$\mathbf{e}$, we have $\widehat{S} = I$ and $\partial_s S = \mathbf{e}$.
More generally, $\widehat{S}$ is an invertible linear operator, and its
adjoint (transpose) satisfies
\begin{equation}
    \bigl\langle v,\; \widehat{S}\,u \bigr\rangle
    = \bigl\langle \widehat{S}^{\top} v,\; u \bigr\rangle
    \qquad \forall\, u,v \in \mathbb{R}^n.
\end{equation}
For a translation, $\widehat{S}^{\top} = \widehat{S}^{-1} = S(\cdot, -s)$, i.e.\ the
inverse shift.  This is the key identity used in the implementation.

\subsection{Adjoint of the composed operator on segment $N$}

The forward map on the last segment is $\widehat{\Phi}_N = \widehat{S} \circ \Phi_N$.
Its transpose is therefore
\begin{equation}
    \bigl(\widehat{S} \circ \Phi_N\bigr)^{\top}
    = \Phi_N^{\top} \circ \widehat{S}^{\top}.
\end{equation}
Consequently, when computing $J(z)^\top \cdot w$, the contribution
$-\widehat{\Phi}_N^{\top} w_N$ is evaluated by \emph{first} applying the inverse
shift $\widehat{S}^{\top}$ to $w_N$ and \emph{then} integrating the adjoint
equations backward through the segment:
\begin{equation}
    -\widehat{\Phi}_N^{\top} w_N
    = -\Phi_N^{\top}\bigl(\widehat{S}^{\top} w_N\bigr).
\end{equation}

\subsection{Full adjoint action $J(z)^\top \cdot w$}

Let $w = (w_1, \dots, w_N,\, w_T,\, w_s)$ where $w_i \in \mathbb{R}^n$ are the
segment components and $w_T, w_s \in \mathbb{R}$ are the phase-condition components.
Transposing~\eqref{eq:jfw_inner}--\eqref{eq:jfw_last} block-wise gives, for each
segment $i = 1,\dots,N$,
\begin{align}
    (J^\top w)_i &=
        -\Phi_i^{\top}\,\widetilde{w}_i \;+\; w_{i-1},
        \qquad\text{with } w_0 := w_N, \\[4pt]
    \widetilde{w}_i &=
        \begin{cases}
            w_i,                 & i < N, \\[2pt]
            \widehat{S}^{\top}\,w_N, & i = N \;\;(\text{NS}=2).
        \end{cases}
\end{align}
The scalar components are
\begin{align}
    (J^\top w)_T &= -\frac{1}{N}\sum_{i=1}^{N-1}
                     \bigl\langle f(\varphi_i),\, w_i \bigr\rangle
                   -\frac{1}{N}\bigl\langle \widehat{f}_N,\, w_N \bigr\rangle, \\[4pt]
    (J^\top w)_s &= \bigl\langle \partial_s S(\varphi_N, s),\; w_N \bigr\rangle
                   \qquad (\text{NS}=2\text{ only}).
\end{align}
The phase-condition rows add the usual contributions
$(J^\top w)_1 \;\;+\!\!=\; f(u_1)\,w_T + g(u_1)\,w_s$.

\subsection{Implementation}

In the code, the adjoint cache (\texttt{AdjointIterSolCache}) computes
$J(z)^\top \cdot w$ matrix-free via the function \texttt{mul!}.
The crucial line
\begin{center}
\texttt{NS == 2 \&\& i == N \&\& S(out[i], -z0.d[2])}
\end{center}
applies the inverse shift $\widehat{S}^{\top} = S(\cdot, -s)$ to the adjoint
variable \emph{before} the backward adjoint integration on segment $N$,
exactly implementing $\Phi_N^{\top} \circ \widehat{S}^{\top}$.
The time-derivative terms $f(\varphi_i)$ (and $\widehat{f}_N$) are pre-computed
in \texttt{dxTdT[i]} during the forward \texttt{update!}, and the derivative
$\partial_s S(\varphi_N,s)$ is evaluated on the fly via the operator
\texttt{D[2]}.


We use the L-BFGS method, a quasi-Newton algorithm that approximates the Hessian
\begin{equation*}
    H(z)=\nabla^2\phi(z) = J(z)^\top J(z) + \sum_{i=1}^{nN+1} F_i(z)\,\nabla^2 F_i(z)
\end{equation*}
using only gradient information. At each iteration, L-BFGS computes a search direction 
\begin{equation*}
    d_k \approx -H_k^{-1}\,\nabla\phi(z_k)
\end{equation*}
where $H_k$ is an approximation of the Hessian $H(z_k)$ using the stored history of at most $m$ pairs $\{(s_i, y_i)\}$ where 
\begin{align*}
    s_i &= z_{i+1}-z_i\\
    y_i &= \nabla\phi(z_{i+1})-\nabla\phi(z_i).
\end{align*}
The gradient is evaluated matrix-free (see Section \ref{sec:discrete_adjoint}) as $\nabla\phi(z) = J(z)^\top F(z)$ using the discrete adjoint. 
A backtracking line search determines the step length $\lambda$ satisfying the Armijo condition, and the state is updated as
\begin{equation*}
    z_{k+1} = z_k + \lambda d_k
\end{equation*}

\subsubsection*{Pseudo-codes}

\begin{algorithm}[ht]
\caption{L-BFGS for $\min_z \frac{1}{2}\|F(z)\|^2$}
\label{alg:lbfgs_opt}
\begin{algorithmic}[1]
\State \textbf{Input:} initial guess $z_0$, memory size $m$, tolerances $\varepsilon_e, \varepsilon_d$
\State Allocate history buffers $\{s_i\}_{i=1}^m, \{y_i\}_{i=1}^m$, set $n_{\text{hist}} \gets 0$
\State Compute $F_0 \gets F(z_0)$, $\nabla\phi_0 \gets J(z_0)^\top F_0$, $e \gets \|F_0\|$
\For{$k = 0,1,2,\dots$ until convergence}
    \If{$e < \varepsilon_e$} \State \textbf{break} \EndIf
    \State Compute search direction $d_k$ via two-loop recursion (Alg.~\ref{alg:two_loop})
    \State $z_{\text{prev}} \gets z_k$, \; $\nabla\phi_{\text{prev}} \gets \nabla\phi_k$
    \State Line search: find $\lambda$ satisfying $\phi(z_k + \lambda d_k) < \phi(z_k)$
    \State $z_{k+1} \gets z_k + \lambda d_k$
    \State Compute $\nabla\phi_{k+1} \gets J(z_{k+1})^\top F(z_{k+1})$, \; $e \gets \|F(z_{k+1})\|$
    \State Update history with $(s = z_{k+1}-z_k,\; y = \nabla\phi_{k+1}-\nabla\phi_k)$
    \If{$\|s\| < 10^{-14}$ \textbf{or} $\|y\| < 10^{-14}$ \textbf{or} $\langle s,y\rangle \le 0$} \State \textbf{skip} \EndIf
    \State $\text{idx} \gets (k \bmod m) + 1$; \; $s_{\text{idx}} \gets s$; \; $y_{\text{idx}} \gets y$; \; $n_{\text{hist}} \gets \min(n_{\text{hist}}+1,\, m)$
    \If{$\|\lambda d_k\| < \varepsilon_d$} \State \textbf{break} \EndIf
\EndFor
\State \textbf{return} $z_k$
\end{algorithmic}
\end{algorithm}

\begin{algorithm}[ht]
\caption{L-BFGS two-loop recursion --- computes $d = -H_k^{-1}\nabla\phi$}
\label{alg:two_loop}
\begin{algorithmic}[1]
\State \textbf{Input:} gradient $\nabla\phi$, history $\{(s_i,y_i)\}_{i=1}^{n}$, step counter $k$, memory $m$
\State $n \gets \min(k,\, m)$ \Comment{effective history length}
\State $q \gets -\nabla\phi$
\For{$i = n,\, n-1,\, \dots,\, 1$} \Comment{first loop (backward)}
    \State $\rho_i \gets 1 / \langle y_i, s_i \rangle$
    \State $\alpha_i \gets \rho_i \langle s_i, q \rangle$
    \State $q \gets q - \alpha_i y_i$
\EndFor
\State $\gamma \gets \begin{cases} \langle s_n, y_n \rangle / \langle y_n, y_n \rangle & \text{if } n > 0 \\ 1 & \text{otherwise} \end{cases}$ \Comment{initial scaling}
\State $q \gets \gamma\, q$
\For{$i = 1,\, 2,\, \dots,\, n$} \Comment{second loop (forward)}
    \State $\beta \gets \rho_i \langle y_i, q \rangle$
    \State $q \gets q + (\alpha_i - \beta) s_i$
\EndFor
\State \textbf{return} $d \gets q$
\end{algorithmic}
\end{algorithm}

Note that in the implementation, the index computation uses $\texttt{idx} = (k-1) \bmod m + 1$ (one-based indexing) to correctly wrap around the circular history buffer, and the history vectors $s_i, y_i$ are stored as \texttt{MVector} objects containing both the $N$ segment states and the scalar period.


\section{Implementation}

\subsection{L-BFGS Method}


\section{Numerical Experiments}

\section{Results}

\section{Outlook}


\end{document}